\documentclass[tikz,border=0pt]{standalone}
\usepackage{ctex}
\usepackage{tikz}

% ===== 自定义金色系 =====
\definecolor{gold}{RGB}{255, 215, 0}
\definecolor{darkgold}{RGB}{184, 134, 11}
\definecolor{deepred}{RGB}{139, 0, 0}
\definecolor{brightgold}{RGB}{255, 240, 150}

% ===== 手绘五角星宏 =====
% #1=x, #2=y, #3=半径, #4=填充色, #5=描边色
\newcommand{\drawstar}[5]{%
	\begin{scope}[shift={({#1},{#2})}]
		\fill[#4, draw=#5, line width=0.5pt]
		(90:#3)
		\foreach \a in {162,234,306,18} {
			-- (\a:#3)
			-- (\a+36:0.382*#3)
		} -- cycle;
	\end{scope}
}

% ===== 四角回纹角花宏 =====
% #1=x, #2=y, #3=旋转角
\newcommand{\cornerornament}[3]{%
	\begin{scope}[shift={(#1,#2)}, rotate=#3]
		\draw[gold, line width=1.2pt]
		(0,0) -- (1.2,0) -- (1.2,0.35) -- (0.35,0.35) -- (0.35,1.2) -- (0,1.2) -- cycle;
		\draw[gold!80!white, line width=0.6pt]
		(0.15,0.15) -- (0.95,0.15) -- (0.95,0.25) -- (0.25,0.25) -- (0.25,0.95) -- (0.15,0.95) -- cycle;
		\fill[gold] (0.6,0.6) circle (0.07);
	\end{scope}
}

\begin{document}

	\begin{tikzpicture}[scale=1]

		% ============================================================
		% ===== 背景层（无白边，铺满画布） =====
		% ============================================================
		\fill[red!90!black] (-0.2,-0.2) rectangle (16.2,24.2);
		\fill[deepred] (-0.2,-0.2) rectangle (16.2,24.2);

		% 中心光晕
		\fill[red!50!black, opacity=0.55] (8,13) circle (13);
		\fill[red!60!black, opacity=0.45] (8,13) circle (10);
		\fill[red!70!black, opacity=0.35] (8,13) circle (7);

		% 底部渐变横带
		\foreach \i in {0,0.5,...,6} {
			\fill[red!75!black, opacity={0.05+0.06*\i}]
			(-0.2,\i-0.2) rectangle (16.2,\i+0.5);
		}

		% 顶部暗角
		\fill[red!95!black, opacity=0.45] (8,24) circle (9);
		\fill[red!98!black, opacity=0.35] (8,24) circle (6);

		% ============================================================
		% ===== 放射光芒（背景层） =====
		% ============================================================
		\foreach \angle in {0,15,...,345} {
			\draw[gold!25!red, line width=0.6pt, opacity=0.35]
			(8,13) -- ++(\angle:18);
		}
		\foreach \angle in {7.5,22.5,...,352.5} {
			\draw[gold!40!red, line width=1.2pt, opacity=0.18]
			(8,13) -- ++(\angle:16);
		}

		% ============================================================
		% ===== 底部扇形光芒 =====
		% ============================================================
		\foreach \angle in {200,210,...,340} {
			\draw[brightgold, line width=0.5pt, opacity=0.25]
			(8,0) -- ++(\angle:12);
		}

		% ============================================================
		% ===== 金色粒子 / 礼花星点 =====
		% ============================================================
		\foreach \x/\y/\r in {
			2.2/22.5/0.06, 4.5/21.8/0.05, 11.5/22.2/0.07, 13.8/21.5/0.05,
			1.5/16.5/0.05, 14.5/16.8/0.06, 2.8/12.5/0.04, 13.2/12.2/0.05,
			1.8/7.5/0.05, 14.2/7.8/0.04, 3.2/9.5/0.04, 12.8/9.2/0.05,
			5.5/18.5/0.05, 10.5/18.2/0.04, 6.8/22.8/0.05, 9.2/22.5/0.06,
			4.2/3.2/0.04, 11.8/3.5/0.05, 2.5/5.2/0.04, 13.5/5.5/0.05
		} {
			\fill[brightgold, opacity=0.85] (\x,\y) circle (\r);
		}

		% ============================================================
		% ===== 多层华丽边框 =====
		% ============================================================
		\draw[gold, line width=5pt] (0.3,0.3) rectangle (15.7,23.7);
		\draw[brightgold, line width=1.5pt] (0.55,0.55) rectangle (15.45,23.45);
		\draw[gold!80!white, line width=0.8pt] (0.75,0.75) rectangle (15.25,23.25);
		\draw[gold!70!white, line width=0.6pt, dashed, dash pattern=on 5pt off 3pt]
		(1.0,1.0) rectangle (15.0,23.0);

		% 内层细金边
		\draw[gold!60!white, line width=0.4pt] (1.2,1.2) rectangle (14.8,22.8);

		% ============================================================
		% ===== 四角回纹角花 =====
		% ============================================================
		\cornerornament{1.0}{1.0}{0}
		\cornerornament{15.0}{1.0}{90}
		\cornerornament{15.0}{23.0}{180}
		\cornerornament{1.0}{23.0}{270}

		% ============================================================
		% ===== 四角宝顶圆环 =====
		% ============================================================
		\foreach \x/\y in {0.75/0.75, 15.25/0.75, 0.75/23.25, 15.25/23.25} {
			\fill[gold] (\x,\y) circle (0.18);
			\draw[gold!80!white, line width=0.6pt] (\x,\y) circle (0.35);
			\draw[gold!60!white, line width=0.4pt] (\x,\y) circle (0.55);
			\draw[brightgold, line width=0.3pt] (\x,\y) circle (0.7);
		}

		% 上下中央宝顶装饰
		\foreach \y in {0.75, 23.25} {
			\fill[gold] (8,\y) circle (0.18);
			\draw[gold!80!white, line width=0.6pt] (8,\y) circle (0.35);
			\drawstar{8}{\y}{0.3}{gold!90!white}{darkgold!80!black}
		}

		% ============================================================
		% ===== 顶部双层卷轴飘带 =====
		% ============================================================
		% 外层阴影
		\fill[black, opacity=0.4]
		(2.4,22.0) -- (13.6,22.0) -- (13.2,21.1) -- (2.8,21.1) -- cycle;
		% 外层飘带
		\fill[gold!85!white]
		(2.3,22.1) -- (13.7,22.1) -- (13.3,21.2) -- (2.7,21.2) -- cycle;
		\fill[gold!65!white]
		(2.3,22.1) -- (1.4,21.65) -- (2.7,21.2) -- cycle;
		\fill[gold!65!white]
		(13.7,22.1) -- (14.6,21.65) -- (13.3,21.2) -- cycle;
		% 内层飘带
		\fill[deepred]
		(3.0,21.85) -- (13.0,21.85) -- (12.7,21.45) -- (3.3,21.45) -- cycle;
		\node[font=\bfseries\zihao{3}, text=gold] at (8,21.65)
		{中国人民解放军建军};

		% 顶部副飘带
		\fill[black, opacity=0.35]
		(4.0,20.6) -- (12.0,20.6) -- (11.7,20.0) -- (4.3,20.0) -- cycle;
		\fill[brightgold]
		(3.9,20.7) -- (12.1,20.7) -- (11.8,20.1) -- (4.2,20.1) -- cycle;
		\node[font=\bfseries\zihao{5}, text=deepred] at (8,20.4)
		{1927 — 2027 · 百年荣光};

		% ============================================================
		% ===== 左右垂幔 =====
		% ============================================================
		% 左侧垂幔
		\fill[gold!70!white, opacity=0.85]
		(1.2,22.8) .. controls (1.8,20) and (1.5,16) .. (1.2,12)
		.. controls (1.5,10) and (1.8,8) .. (1.2,6)
		.. controls (1.5,4) and (1.8,3) .. (1.2,2.2)
		-- (0.6,2.2) .. controls (0.3,3) and (0.6,4) .. (0.6,6)
		.. controls (0.3,8) and (0.6,10) .. (0.6,12)
		.. controls (0.3,16) and (0.6,20) .. (0.6,22.8) -- cycle;
		% 右侧垂幔
		\fill[gold!70!white, opacity=0.85]
		(14.8,22.8) .. controls (14.2,20) and (14.5,16) .. (14.8,12)
		.. controls (14.5,10) and (14.2,8) .. (14.8,6)
		.. controls (14.5,4) and (14.2,3) .. (14.8,2.2)
		-- (15.4,2.2) .. controls (15.7,3) and (15.4,4) .. (15.4,6)
		.. controls (15.7,8) and (15.4,10) .. (15.4,12)
		.. controls (15.7,16) and (15.4,20) .. (15.4,22.8) -- cycle;

		% ============================================================
		% ===== 中心军徽 =====
		% ============================================================
		\draw[gold, line width=3pt] (8,15.5) circle (3.7);
		\draw[brightgold, line width=1pt] (8,15.5) circle (3.95);
		\draw[gold!80!white, line width=0.8pt] (8,15.5) circle (4.15);
		\draw[gold!60!white, line width=0.5pt, dashed, dash pattern=on 3pt off 2pt]
		(8,15.5) circle (4.4);

		% 大五角星阴影
		\drawstar{8.08}{15.42}{2.2}{black!55}{black!55}
		% 大五角星主体
		\drawstar{8}{15.5}{2.2}{gold}{darkgold}
		% 高光小星
		\drawstar{8}{15.5}{1.6}{brightgold!70!gold}{darkgold}

		% 内嵌红色圆
		\fill[deepred] (8,15.5) circle (1.05);
		\draw[brightgold, line width=0.5pt] (8,15.5) circle (1.05);

		% 八一文字
		\node[font=\bfseries\zihao{2}, text=gold] at (8,15.5) {八一};

		% 环绕五颗小星
		\foreach \angle in {90,162,234,306,18} {
			\pgfmathsetmacro{\sx}{8+3.1*cos(\angle)}
			\pgfmathsetmacro{\sy}{15.5+3.1*sin(\angle)}
			\drawstar{\sx}{\sy}{0.55}{gold!90!white}{darkgold!80!black}
		}

		% ============================================================
		% ===== 左右竖排对联式标语 =====
		% ============================================================
		\node[font=\bfseries\zihao{3}, text=gold, rotate=90] at (2.6,13.5)
		{铁血铸军魂};
		\node[font=\bfseries\zihao{3}, text=gold, rotate=-90] at (13.4,13.5)
		{丹心照汗青};

		% ============================================================
		% ===== 主标题区 =====
		% ============================================================
		% 标题外描边
		\node[font=\bfseries\zihao{0}, text=darkgold, opacity=0.9,
		xshift=3pt, yshift=-3pt] at (8,11.2) {建军百年};
		% 标题阴影
		\node[font=\bfseries\zihao{0}, text=black!55, opacity=0.6,
		xshift=1.5pt, yshift=-1.5pt] at (8,11.2) {建军百年};
		% 标题主体
		\node[font=\bfseries\zihao{0}, text=gold] at (8,11.2) {建军百年};

		\node[font=\bfseries\zihao{2}, text=brightgold] at (8,9.6)
		{1927 \; — \; 2027};

		% ============================================================
		% ===== 分隔装饰线 =====
		% ============================================================
		\draw[gold, line width=2pt] (3.5,8.8) -- (12.5,8.8);
		\draw[brightgold, line width=0.8pt] (4.5,8.55) -- (11.5,8.55);
		\drawstar{8}{8.8}{0.4}{gold!90!white}{darkgold!80!black}
		\drawstar{4.8}{8.8}{0.25}{gold!80!white}{darkgold!70!black}
		\drawstar{11.2}{8.8}{0.25}{gold!80!white}{darkgold!70!black}

		% ============================================================
		% ===== 副标题 =====
		% ============================================================
		\node[font=\bfseries\zihao{3}, text=gold!95!white] at (8,7.6)
		{峥嵘岁月 \; · \; 铁血军魂};
		\node[font=\zihao{4}, text=gold!85!white] at (8,6.8)
		{传承红色基因 \; · \; 担当强军重任};

		% ============================================================
		% ===== 底部三大标语 =====
		% ============================================================
		\foreach \x/\txt in {3.5/听党指挥, 8/能打胜仗, 12.5/作风优良} {
			\draw[gold!80!white, line width=1.2pt, rounded corners=4pt]
			(\x-1.7,4.2) rectangle (\x+1.7,5.6);
			\draw[brightgold, line width=0.5pt, rounded corners=3pt]
			(\x-1.55,4.35) rectangle (\x+1.55,5.45);
			\node[font=\bfseries\zihao{4}, text=gold] at (\x,4.9) {\txt};
		}

		% ============================================================
		% ===== 底部长城剪影 =====
		% ============================================================
		\begin{scope}[shift={(0,0)}]
			\fill[gold!55!red, opacity=0.55]
			(2.0,2.8) -- (2.0,3.3) -- (2.4,3.3) -- (2.4,3.0)
			-- (3.0,3.0) -- (3.0,3.5) -- (3.4,3.5) -- (3.4,3.2)
			-- (4.0,3.2) -- (4.0,3.7) -- (4.4,3.7) -- (4.4,3.3)
			-- (5.0,3.3) -- (5.0,3.8) -- (5.4,3.8) -- (5.4,3.4)
			-- (6.0,3.4) -- (6.0,3.9) -- (6.4,3.9) -- (6.4,3.5)
			-- (7.0,3.5) -- (7.0,4.0) -- (7.4,4.0) -- (7.4,3.6)
			-- (8.0,3.6) -- (8.0,4.1) -- (8.4,4.1) -- (8.4,3.6)
			-- (9.0,3.6) -- (9.0,4.0) -- (9.4,4.0) -- (9.4,3.5)
			-- (10.0,3.5) -- (10.0,3.9) -- (10.4,3.9) -- (10.4,3.4)
			-- (11.0,3.4) -- (11.0,3.8) -- (11.4,3.8) -- (11.4,3.3)
			-- (12.0,3.3) -- (12.0,3.7) -- (12.4,3.7) -- (12.4,3.2)
			-- (13.0,3.2) -- (13.0,3.5) -- (13.4,3.5) -- (13.4,3.0)
			-- (14.0,3.0) -- (14.0,2.8) -- cycle;
		\end{scope}

		% ============================================================
		% ===== 底部年份与落款 =====
		% ============================================================
		\node[font=\bfseries\zihao{3}, text=brightgold] at (8,2.3)
		{1927 \; — \; 2027};

		\node[font=\zihao{5}, text=gold!70!white] at (8,1.0)
		{纪念中国人民解放军建军一百周年};

		% ============================================================
		% ===== 底部小星装饰 =====
		% ============================================================
		\foreach \x in {3.5,4.7,5.9,10.1,11.3,12.5} {
			\drawstar{\x}{0.5}{0.2}{gold!90!white}{darkgold!80!black}
		}

		% ============================================================
		% ===== 两侧麦穗勋章装饰 =====
		% ============================================================
		\begin{scope}[shift={(2.0,11)}, rotate=0]
			\draw[gold!85!white, line width=2pt]
			(0,4.5) .. controls (-0.7,2.5) and (-0.5,0.5) .. (0,-1.5);
			\foreach \y in {4.0, 3.4, 2.8, 2.2, 1.6, 1.0, 0.4, -0.2, -0.8} {
				\draw[gold!80!white, line width=0.7pt]
				(0,\y) .. controls (-0.55,\y+0.4) and (-1.0,\y+0.3) .. (-1.3,\y+0.7);
				\fill[brightgold] (-1.25,\y+0.65) circle (0.08);
			}
		\end{scope}

		\begin{scope}[shift={(14.0,11)}, rotate=180]
			\draw[gold!85!white, line width=2pt]
			(0,4.5) .. controls (-0.7,2.5) and (-0.5,0.5) .. (0,-1.5);
			\foreach \y in {4.0, 3.4, 2.8, 2.2, 1.6, 1.0, 0.4, -0.2, -0.8} {
				\draw[gold!80!white, line width=0.7pt]
				(0,\y) .. controls (-0.55,\y+0.4) and (-1.0,\y+0.3) .. (-1.3,\y+0.7);
				\fill[brightgold] (-1.25,\y+0.65) circle (0.08);
			}
		\end{scope}

		% ============================================================
		% ===== 勋章圆环（上下对称） =====
		% ============================================================
		\foreach \y in {18.8, 5.8} {
			\draw[gold, line width=1.2pt] (8,\y) circle (0.55);
			\draw[brightgold, line width=0.5pt] (8,\y) circle (0.4);
			\drawstar{8}{\y}{0.28}{gold!90!white}{darkgold!80!black}
		}

		% ============================================================
		% ===== 飘带装饰 =====
		% ============================================================
		\draw[gold!70!white, line width=0.8pt, opacity=0.8]
		(2.5,19.5) .. controls (5,19.1) and (11,19.1) .. (13.5,19.5);
		\draw[brightgold, line width=0.5pt, opacity=0.7]
		(2.8,19.1) .. controls (5,18.8) and (11,18.8) .. (13.2,19.1);

	\end{tikzpicture}

\end{document}